\section{Ordered labels and stability of selected rays}\label{sec:labels}
All uses of AH Sections 6--10 in this section are \emph{proof-level reuse}.
AH's standing A1/A2 assumptions include finite-energy conditions not satisfied
by $s$. Section~\ref{sec:uses} lists each use and its replacement premise;
we do not apply its final structural theorem as a black box.

The reference families, uniqueness and planar order give AH Section 6.1's
cyclic ray order. Sections 6.2--6.4 average source flows over independent
auxiliary Voronoi cells. The product structure is the cap law times those
auxiliary uniforms, not independence of capped edges. The conditional-mean
construction in Sections 6.5--6.7 gives monotone labels consistent under
coalescence. Section 7 uses planar trapping, ergodic density and extended
shape to control noncrossing selectors. These arguments introduce no edge
resampling. The one downward operation occurs in AH 8.6.

\subsection{The replacement for AH 8.6}
For an edge $e$, independently resample $t_e$ conditional on its open/closed
status, leaving all $V$ fixed. This kernel is reversible for the underlying
product law. Every status, hence every $L_f$ and $K_f$, stays fixed; only
$s_e$ can change. If $e$ is usable and uncapped, its cap decreases by more
than $\delta=g/2$ with probability at least $p_*$, by sampling its low band.
The fixed $\delta,p_*$ are chosen before the root and finite witness paths
in AH's bad event.

Let $\mathcal F_n$ be the originals and auxiliaries on edges with both
endpoints in $Q_n$. For annular edges outside this set, the kernel preserves
$\mathcal F_n$ and is reversible conditionally on it. Let
$A=A(v,P_1,P_2,P_3,\delta)$ be AH's finite-witness event, with its witness
paths inside $Q_n$. Its deterministic Lemmas 8.4--8.5 use exactly one
decreased edge: a successful decrement on the relevant extremal ray $g_1$
destroys $A$ and places that edge on the new clockwise extremal. They
therefore apply to this exact single-cap-edge change.

Let $C_1(n)$ be the uniform length event of Lemma~\ref{lem:metric}, and let
$C'_1(n)$ use $3Ln$ in place of $Ln$. Lowering one edge implies
$C_1\Rightarrow C'_1$: a new geodesic using the edge splits into two
old-minimal pieces with endpoints in $Q_{2n}$, so has at most $2Ln+1$
edges. An avoiding path is already old-minimal. Usable-edge counting in
the original iid field gives at least half usable edges on every long
simple path from the box, with probability tending to one. All $v$-rays
are uncapped beyond finite $E_v$. Their crossings of the annulus thus have
at least $\gamma n$ eligible edges, say $\gamma=1/3$, on an event $C_2(n,v)$
of probability tending to one.

Writing $D_e$ for eligibility, conditional reversibility gives
\[
 p_*\PP(A\cap C_1\cap D_e\cap\{e\in g_1\}\mid\mathcal F_n)
 \le\PP(A^c\cap C'_1\cap\{e\in g_1\}\mid\mathcal F_n).
\]
Sum over annular edges. On $C'_1$ the part of $g_1$ between its first and
last visits to the outer box has at most $3Ln$ edges. Consequently
\begin{equation}\label{eq:lowering}
 \PP(A\mid\mathcal F_n)\le
 \PP(C_1^c\mid\mathcal F_n)+\PP(C_2^c\mid\mathcal F_n)
 +\frac{3L}{p_*\gamma}\PP(A^c\mid\mathcal F_n).
\end{equation}
The two error terms tend to zero in probability, whereas martingale
convergence makes the conditional probability of $A$ tend to $\ind_A$.
If $\PP(A)>0$, \eqref{eq:lowering} is impossible on a positive-probability
part of $A$ for large $n$. Thus $\PP(A)=0$. Countability of the finite
witnesses yields AH 8.2 with its original simultaneous quantifiers.

The remaining proofs of Sections 8.2--8.4 now use order, compact limits and
countability to give a deterministic closed label set $\Lcal_*$, the two
extremal selectors at each fixed label, and their continuity properties.
The proof of AH 9.2 constructs a stationary increment/path extension.
Lemma~\ref{lem:cocycle}, rather than an assertion of extension ergodicity,
provides its linearity. Its supporting-arc argument and deterministic
selector direction set yield one of the eight eventual half-planes.

\subsection{S: common prediction sites and symmetric label transfer}
Consider the original upward coupling of Lemma~\ref{lem:lift}. It is
generally not stationary as a joint law, but both marginals are absolutely
continuous with respect to the original law. Its cap difference set is finite.
Reference stability from Section~\ref{sec:reference}, together with the
geometric proof of AH 9.7 and extended metric/Busemann bounds, gives a
translate of a cone of positive deterministic angle in which the two
bounding reference rays agree and the sector contains no changed edge.
For agreement, the old and new reference rays avoid the change set and
share a tail; their prefixes to a common tail vertex have the same costs
in both fields, so finite uniqueness identifies them. This replaces the
specific AH 9.6 conclusion used in the proof of 9.7.

Fix a label and the finite neighboring-selector list used in AH 9.8.
Choose a finite stub length so neighboring unequal labels have different
stubs with probability at least $1-\eta$, then a finite cap-patch radius
whose maximum-likelihood predictions of those stubs are simultaneously
correct with probability at least $1-\eta$, adjusting $\eta$ for the
finite list. Martingale convergence supplies this prediction property;
fixed tie-breaking makes the predictors translation-covariant.
Ergodicity gives bad-site density at most $\eta$ in each marginal on a
full event. Absolute continuity transfers that full event to both fields.
Their joint bad-site union therefore has upper density at most $2\eta$,
even though the coupling is not stationary. Choose $2\eta$ below the cone's
density, and discard the finitely many patches meeting the changed set.
Some cone site is good in both fields and has identical patches, hence
identical actual selector stubs.

The planar comparison in the proof of AH 9.8 now traps new extremal labels
between their prescribed old neighbors. New realized labels belong to
$\Lcal_*$ by absolute continuity of the new marginal. If $\alpha$ is
approached from below, use a countable sequence of lower labels. If it has
an immediate predecessor $\beta$, the strict comparison ``new label
$>\beta$'' and membership in $\Lcal_*$ imply new label $\ge\alpha$.
Use the symmetric successor convention above; reference-equality cases
are handled directly. This avoids a nonexistent approximating sequence
at an isolated label. Monotone squeezing transfers every intervening
$\alpha$-ray, not only the two extremals.

The premises of this label comparison are symmetric: both marginals are
absolutely continuous, the cap difference is finite, and corresponding
reference rays coalesce. Apply it in both orientations. The reverse-label
statement does not require a reverse upward resampling kernel.

\begin{lemma}[Finite-change tails]\label{lem:tails}
If two positive coercive metrics $w,w'$ agree off finite $F_0$, every
$w$-ray has a $w'$-geodesic tail, and conversely. The cutoff is uniform
over rays from a fixed root.
\end{lemma}
\begin{proof}
If $F_0$ is empty the assertion is immediate. Otherwise let $V_0$ contain the endpoints of $F_0$, let
$\Delta=\sum_{F_0}|w'_e-w_e|$, and put
$C=\max_{v\in V_0}T_w(x,v)$. A simple competitor from a late $a$ to a later
$b$ on an $x$-ray that meets $V_0$ has new cost at least
\[
 T_w(x,a)+T_w(x,b)-2C-\Delta
 =T_w(a,b)+2T_w(x,a)-2C-\Delta.
\]
Use the first and last contacts and triangle inequalities for this bound.
Once $T_w(x,a)>C+\Delta/2$ and the tail avoids $V_0$, it cannot improve the
unchanged ray segment. Avoiding competitors have identical costs. Positive
weights permit loop erasure before the one-copy $\Delta$ estimate.
Coercivity supplies the uniform late cutoff; reverse the metrics for the converse.
\end{proof}

\subsection{Entire extremal selectors and the second wall application}
Suppose an old extremal $\alpha$-ray from $v$ is wholly unchanged under
$s'\ge s$. It is a new ray by monotonicity. If the new extremal differed
on the relevant side, Lemma~\ref{lem:tails} and the reverse label comparison
would make its tail an old $\alpha$-ray. The proof of AH 9.9 then says old
finite geodesics from $v$ to late vertices of that tail converge to the old
extremal. Put $v$ and the cap changes in $Q_R$, and take $Q_S$ containing
all old geodesics between vertices of $Q_R$. Once such a finite geodesic
agrees with the old extremal through its first exit from $Q_S$, it is wholly
unchanged: a later return to $Q_R$ is impossible. Finite uniqueness makes
it strictly cheaper than the different new-extremal prefix in $s$, and its
unchanged cost is still strictly cheaper in $s'$, a contradiction.
Thus the entire selected ray is preserved. Prefix convergence alone would
not have sufficed without this no-return argument.

Apply the wall of Section~\ref{sec:wall} to late uncapped extremal triples.
The selector premise is now proved by the preceding paragraph, rather than
left as a historical condition. AH 9.11's mass transport gives coalescence.
The remaining proofs of 9.1 and 9.12--9.13 use deterministic labels,
stationarity and disintegration, without a new energy operation.

\subsection{The all-ray output needed by S}\label{sec:allray}
The proof-level chain in AH Section 10, detailed in Section~\ref{sec:uses},
now supplies a deterministic closed set $\Ccal$ of linear functionals
dominated by $\mu$. On one event every cap ray $p$ has a Busemann function
$B_p$ with
\[
 B_p(0,x)=\rho_p(x)+o(|x|),\qquad\rho_p\in\Ccal,
\]
uniformly at infinity; its limiting directions lie in the contact arc
$\{u:\rho_p(u)=\mu(u)\}$. The uniform all-ray cone bounds come from
countably many finite reference nets, not an uncountable intersection of
fixed-ray assertions. For each \emph{fixed} functional there is a unique
coalescing, backwards-finite family; its null set may depend on that
functional. Backwards finiteness uses the mass-transport proof of AH 4.5.
The presence of every tangent functional uses the explicitly quantified
construction in Section~\ref{sec:reference}. All these statements concern
$s$ and its Busemann functions, not uniform asymptotics for the original $t$.
